\cotefile{157-1}
\scanpage{66}
\mspage{72}

\uline{Th. 3} Soit
\diagram{\begin{tikzcd}[ampersand replacement=\&]
X \arrow[d,"f"'] \& X' \arrow[l,"q"'] \arrow[d,"f'"]\\
Y \& Y' \arrow[l,"p"]
\end{tikzcd}}
cartésien dans $\mathrm{Diag}$. Alors le morph. de chgt. de base
\[
p^* f_* \longrightarrow f'_* q^*
\]
est un iso si $f$ est $\mathbb D$-propre, ou \deleted{\unsure{$f$|$\mathbb D$}} \unsure{$p$|$f$} $\mathbb D$-lisse.

Le cas $\mathbb D$-propre déjà \unsure{traité}. Le cas $\mathbb D$-lisse : \deleted{\unsure{est pro}} \inserted{\unsure{la|sa}} dém. p.~58 est OK.

\figure{\begin{tikzpicture}\draw (0,.32) .. controls (.05,.18) and (.19,.11) .. (.22,0);\end{tikzpicture}}

Comparaison de $X$ et $X^\circ$.

Soit $X$ \deleted{\unsure{petite}} catégorie, on \unsure{va|se} définir une catégorie $\unsure{\mathscr D|\mathscr P}(X)$ et des foncteurs
\diagram{\begin{tikzcd}[ampersand replacement=\&,column sep=small]
\& \& \mathscr D(X) \arrow[dl,"{p_X}"'{name=pl}] \arrow[dr,"{\deleted{\unsure{p}} \unsure{q'_X|p'_X}}"{name=pr}] \& \& \\
\substack{\text{Cat-cofibrant}\\(\unsure{\text{donc}}\ \inserted{\unsure{u\text{-}}}\text{propre})}
\& X \& \& X^\circ \& \substack{\text{Cat-fibrant}\\(\unsure{\text{donc}}\ \text{lisse})}
\arrow[from=pl,to=2-1]\arrow[from=pr,to=2-5,dash]
\end{tikzcd}}
qui seront des \unsure{hot}-équivalences, et des $\mathbb D$-équivalences si $\mathbb D$ est un dérivateur dont le domaine contient $X$ et $\mathscr D(X)$. On pose
\[
\mathrm{Ob}\,\mathscr D(X)=\mathrm{Fl}(X),
\]
donc \unsure{$=$|mêmes} objets que $\underline{\mathrm{Fl}}(X)=\underline{\mathrm{Hom}}(\Delta^1,X)$ ; mais les hom sont différents :
% archive sheet number: 65
\scanpage{67}
\mspage{73}
\diagram{\begin{tikzcd}[ampersand replacement=\&]
y_1 \arrow[r,"u"] \& y_0 \arrow[d,"v_0"]\\
y'_1 \arrow[u,"v_1"] \arrow[r,"u'"] \& y'_0
\end{tikzcd}}

un hom de $u$ dans $u'$ est un couple $\left\{\begin{aligned}v_0&;y_0\to y'_0,\\v_1&:y'_1\to y_1\end{aligned}\right.$ tel que $v_0 u v_1=u'$.

Ainsi source et but $s(u)$ et $b(u)$ sont l'une contravariante, l'autre covariante en $u$. On prend pour $q$ le foncteur source, pour $p$ le foncteur but.

Si on applique ceci à $X^\circ$, on trouve
\diagram{\begin{tikzcd}[ampersand replacement=\&]
\& {\unsure{\mathscr D|\mathscr P}(X^\circ)\unsure{-}} \arrow[dl,"p_{X^\circ}"'] \arrow[dr,"q_{X^\circ}"] \& \\
X^\circ \& \& (X^\circ)^\circ=X
\end{tikzcd}}

On aura
\[
\mathrm{Ob}\,\deleted{\unsure{\mathrm{Fl}}}\mathscr D(X^\circ)=\mathrm{Fl}(X^\circ)\simeq\mathrm{Fl}(X)
\]
\diagram{\begin{tikzcd}[ampersand replacement=\&]
y_1 \arrow[r,"u"] \& y_0 \arrow[d,"v_1",very thick] \& \text{dans }X\\
y'_1 \arrow[u,"v_0"] \arrow[r,"u'"] \& y'_0 \& \text{dans }X
\end{tikzcd}}
et $\unsure{=^{\mathrm s}|=}\,$ flèches au sens de $\mathscr D(X)$ et au sens de $\mathscr D(X^\circ)$.

Donc
\[
\left\{\begin{aligned}
\mathscr D(\unsure{X^\circ|X})&=\mathscr D(X^\circ)\\
p_{X^\circ}&=q_X\\
q_{X^\circ}&=p_X
\end{aligned}\right.,
\]
le passage de $X$ à $X^\circ$ revient à échanger $p_X$ et $q_X$.

La fibre de $\mathscr D(X)$ en $x\in\mathrm{Ob}\,X$ est $\simeq(X_{/x})^\circ$, celle en $y^\circ\in\mathrm{Ob}\,X^\circ$ est $\simeq {}_{y\backslash}X$
\[
\left\{\begin{aligned}
p_X^{-1}(\{x\})&\simeq(X_{/x})^\circ\\
\deleted{\unsure{p}}q_X^{-1}(\{y^\circ\})&\simeq {}_{y\backslash}X_{\deleted{\unsure{/y}}}
\end{aligned}\right.
\]

\margin{NB : $p_X$ est $W$-lisse à fibres $W$-asph, donc dans $W^{\unsure{u|w}}$. $q_X$ est $W$-propre à fibres \unsure{$W$-asph} donc dans $W^u$.}
% archive sheet number: 66
\scanpage{68}
\mspage{74}

On vérifie
\[
\left\{\begin{aligned}
\unsure{\mathscr D|\mathscr P}(X)_{/x}&\simeq\mathscr D(X_{/x}) & x&\in\mathrm{Ob}\,X\\
\mathscr D(X)_{/x^\circ}&\simeq\mathscr D({}_{x\backslash}X) & x^\circ&\in\mathrm{Ob}\,X^\circ
\end{aligned}\right.
\]

Je dis que $p_X$, \inserted{$q_X$} \uline{sont des $\mathbb D$-équiv. strictes}, \inserted{et costrictes} \deleted{\unsure{$q_X$ une} $\mathbb D$-équivalence costricte}. Par les formules précédentes, \deleted{\unsure{on est ramené}} cela signifie que les $\mathscr D(X_{/x})$, $\mathscr D({}_{x\backslash}X)$ sont $\deleted{\unsure{d}}\mathbb D$-contractiles. Donc que si $Z$ a un objet final, $\mathscr D(Z)$ est $\mathbb D$-contractile.

Mais considérons
\diagram{\begin{tikzcd}[ampersand replacement=\&]
\mathscr D(Z) \arrow[dr,"q_Z"] \& \\
\& Z^\circ,
\end{tikzcd}}
il suffit \deleted{\unsure{que}} \unsure{de} prouver que c'est une \inserted{$\mathbb D$-}équiv. stricte (puisque $Z^\circ$, ayant un objet initial, est $\mathbb D$-contractile), donc que les $ {}_{z_{\unsure{0|\circ}}\backslash}\mathscr D(Z)$ sont $\mathbb D$-contractiles

On a par la formule précédente
\[
{}_{z_0\backslash}\mathscr D(Z)\simeq\mathscr D({}_{z_0\backslash}Z)=\mathscr D(T),
\]
où cette fois $T$ a un objet initial \inserted{$\unsure{o|0}$} et \deleted{final.} un objet final $e$. Considérons alors l'objet
\diagram{\begin{tikzcd}[ampersand replacement=\&]o\arrow[r]\& e\end{tikzcd}}
de $\mathscr D(T)$, c'est un objet final, car \deleted{\unsure{$\forall$}} pour tout $x\longrightarrow y$ de $\mathscr D(T)$, il y a une unique flèche
\diagram{\begin{tikzcd}[ampersand replacement=\&]
x\arrow[r] \& y\arrow[d]\\
o\arrow[u]\arrow[r] \& e
\end{tikzcd}}
% archive sheet number: 67
\scanpage{69}
\mspage{75}

Pour résumer

% slanted
\margin{i.e. Hot-asphériques}

\uline{Proposition 1} : $p_X$ et $q_X$ sont des $\deleted{\unsure{\mathbb D}}\mathrm{Hot}$-équivalences strictes, à fibres les $(X_{/x})^\circ$ et $ {}_{\deleted{\unsure{y}}x\backslash}X$ respectivement (donc à fibres contractiles)

\margin{(\unsure{donc aussi|Cela en sera} $\mathbb D$-équiv. strictes si $\unsure{\mathscr D|\mathscr P}(X)$ dans dom de $\mathbb D$)}

\deleted{\uline{Corollaire} On a une équivalence} \deleted{On voudrait}

\uline{Corollaire 1} Pour que $X$ soit $\mathbb D$-contractile, il f. et s. que $X^\circ$ le soit.

Plus généralement

\uline{Corollaire 2} Soit $f:X\longrightarrow Y$. Pour que $f$ soit une $\mathbb D$-équivalence, il f. et s. que $f^\circ$ le soit.

On utilise que le diagramme
\diagram{\begin{tikzcd}[ampersand replacement=\&]
\& \mathscr D(X) \arrow[dl,"p_X"'] \arrow[dr,"q_X"] \& \\
X \& \& X^\circ
\end{tikzcd}}
est fonctoriel en $X$. \unsure{On a donc|Ceci a donc}
\diagram{\begin{tikzcd}[ampersand replacement=\&,column sep=small]
\& \mathscr D(X) \arrow[rr,leftrightarrow,"{\mathscr D(f)}"] \arrow[dl,"p_X"'] \arrow[dd,"q_X"] \& \& \mathscr D(Y) \arrow[dl,"p_Y"] \arrow[dd,"q_Y"]\\
X \arrow[rr,"f",crossing over] \& \& Y \& \\
\& X^\circ \arrow[rr,"f^\circ"] \& \& Y^\circ
\end{tikzcd}}
les deux carrés commutatifs. Comme $p_X,p_Y$ sont des $\mathbb D$-équiv., $f$ est une $\mathbb D$-équiv. sss $\mathscr D(f)$ l'est, et de $=$ pour $f^\circ$.

Ces raisonnements supposent que si $X$ est dans le domaine de $\mathbb D$, $\mathscr D(X)$ aussi. C'est le cas si $\mathrm{Diag}$ est formé des cat. ordonnées finies.

\figure{\begin{tikzpicture}\draw (0,.35) -- (.17,.18) -- (.34,0);\draw[line width=.8pt] (.14,.23) -- (.21,.13);\end{tikzpicture}}
% archive sheet number: 68
\scanpage{70}
\mspage{76}

\uline{$\mathbb D$-fibrations}

\uline{Prop 1.} Soit $f:X\longrightarrow Y$ dans $\mathrm{Diag}$, Conditions équivalentes :

a) Pour \deleted{\unsure{$y\longrightarrow y'$}} tout $y\in Y$, l'inclusion $X_y\longrightarrow X_{/y}$ est une $\mathbb D$-équivalence, et pour toute flèche $y_1^{\deleted{\unsure{*}}}\longrightarrow y_0$ de $Y$, $X_{/y_1}\longrightarrow X_{/y_0}$ est une $\mathbb D$-équivalence

b) Pour tout \deleted{\unsure{faisceau}} coefficient \deleted{\unsure{loc.}} \uline{constant} $F$ sur $Y$, $f_*f^*F$ est loc. constant\textsuperscript{+}, et ``se calcule fibre par fibre'', l'hom can
\[
f_* f^*(F)_y \unsure{\simeq|\xrightarrow{\sim}} H^{\unsure{\bullet|\cdot|\prime}}_{\mathbb D}(X_y,F\mid X_y).
\]

% slanted
\margin{\textsuperscript{+} Par là on entend que \deleted{\unsure{$\forall$}} $u:y\longrightarrow y'$ dans $Y$, $\unsure{C|G}_{y'}\longrightarrow C_y$ est iso. Sans axiome suppl. sur $\mathbb D$, \deleted{\unsure{$C$}} il n'est pas clair que cela implique que les ``restrictions'' de $C$ aux $Y/y$, ou aux $y\backslash Y$, soient des coefficients constants. Voir cependant l'axiome $\unsure{\mathrm{Der}\,3\,a')|\mathrm{Der}\,3\,a)}$, page \unsure{97|47}, qui assure que c'est pareil.}

Ces conditions sont équivalentes, et \unsure{impliquent|implique} \deleted{que} la validité de b) pour \inserted{tout} $F$ \uline{loc.} constant sur $Y$.

La \inserted{première} condition \inserted{du} a) signifie que \inserted{$\forall y\in Y$, et} $\forall M\in\mathrm{Ob}\,A$, \deleted{\unsure{en tant}}
\[
H^\bullet_{\mathbb D}(X_{/y},M)\longrightarrow H^\bullet_{\mathbb D}(X_y,M)
\]
est un iso. Or cette flèche s'interprète comme
\[
f_*(M_X)_{\unsure{x|y}}\longrightarrow H^\bullet_{\mathbb D}(X_y,M)
\]
est iso, i.e. que $f_*(M_X)$ se calcule fibre par fibre. Comme $M_X\simeq f^*(M_Y)$, c'est la deuxième condition b).

La deuxième condition du a), pour $y_1\longrightarrow y_0$ signifie que pour $M$ dans $A$,
\[
H^\bullet_{\mathbb D}(X_{/y_0},M)\longrightarrow H^\bullet_{\mathbb D}(X_{/y_1},M)
\]
est iso, \deleted{donc que} or cette flèche s'interprète comme
\[
f_*(M_X)_{y_0}\longrightarrow f_*(M_X)_{y_1},
\]
que ce soit \inserted{tjrs} un iso. est la première condition b). Que b) soit valable aussi pour $F$ loc. constant sur $Y$ est immédiat, car c'est une condition de nature locale p.r. à $Y$.
% archive sheet number: 69
\scanpage{71}
\mspage{77}

La proposition duale fait appel aux foncteurs can
\[
{}_{y_0\backslash}X\longrightarrow {}_{y_1\backslash}X
\qquad\text{et}\qquad X_y\longrightarrow {}_{y\backslash}X,
\]
et aux faisceaux $f_! f^*M$, pour $M$ constant ou \unsure{loc|loc.} constant sur $Y$.

\uline{Définition 1} Si les conditions équiv. de la prop. 1 sont satisfaites, on dit que $f$ est une $\mathbb D$-fibration cohomologique \deleted{\unsure{faible}} \inserted{(ou $\mathbb D$-fibration \uline{gauche})}. Si les conditions duales sont satisfaites, on dit que c'est une $\mathbb D$-fibration homologique (ou $\mathbb D$-fibration \uline{droite}).

\uline{NB} \unsure{Les} $\mathbb D$-fibrations cohomo\uline{lo}giques sont stables par tt chgt de base qui \deleted{\unsure{est}} un iso local. Les $\mathbb D$-fibrations homologiques par \deleted{\unsure{tout}} chgt de base qui est un co-iso local. Mais ce n'\deleted{\unsure{est pas une}} sont pas des notions stables par \uline{tout} changt de base. Mais \unsure{la} notion de $\mathbb D$-fibration coh. est stable par changt de base $\mathbb D$-lisse, et celle de $\mathbb D$-fibration homologique par chgt de base $\mathbb D$-propre.

\margin{? donner exemples !}

Je vais essayer de construire un exemple où le chgt de base coule, en prenant $Y=\Delta_1$. Se donner $X$ sur $\Delta_1$. revient à se donner une sous-catégorie ouverte $X_0$, et sa complémentaire fermée $X_1$
\diagram{\begin{tikzcd}[ampersand replacement=\&,column sep=small]
\& X \& \\
X_1 \arrow[ur,phantom,"{\unsure{\subset|-}}" description,sloped] \arrow[dd,dotted,dash] \& \& X_0 \arrow[ul,phantom,"{\unsure{\supset|-}}" description,sloped] \arrow[dd,dashed,dash]\\
\& \& \\
1 \& \& 0 \arrow[ll]
\end{tikzcd}}
Dire que $X\longrightarrow Y=\Delta_1$ est \deleted{coh.} $\mathbb D$-fibr. coh. signifie simplement
% archive sheet number: 70
\scanpage{72}
\mspage{78}

que les inclusions

\margin{$\unsure{*|x}$}
\[
X_0\xrightarrow{\unsure{j|\partial}}X,\quad X_1\xrightarrow{i}X
\]
sont des $\mathbb D$-équivalences. Prenons un chgt de base $Y'\unsure{\longleftrightarrow|\longrightarrow}Y=\Delta^1$, i.e un $Y'$ avec $Y'_0,Y'_1$
\diagram{\begin{tikzcd}[ampersand replacement=\&]
Y'_1 \arrow[d,dashed,dash] \& Y' \arrow[l,"{\unsure{i'|\partial'}}"'] \arrow[r,"j'"] \& Y'_0 \arrow[d,dash]\\
1 \& \& 0 \arrow[ll]
\end{tikzcd}}

On voit que la condition $Y'$ lisse sur $\Delta_1$ signifie que $\forall y'_1\in\mathrm{Ob}\,Y'_1$, la catégorie $Y'_0/y'_1$ des objets $y_0$ \uline{de $Y'_0$} au dessus de $y'_1$ (i.e munis de $y_0\longrightarrow y'_1$) est contractile.

\margin{$\longrightarrow$}

\uline{Cela exige} que si $Y'_1\ne\varnothing$, \deleted{\unsure{$Y'_0$}} alors $Y'_0\ne\varnothing$ \unsure{---|$=$} mais un exemple avec $Y'_0=\varnothing$ \inserted{ou $Y'_1=\varnothing$} n'est pas intéressant (\unsure{car} \deleted{\unsure{tout $X'\longrightarrow$}} une projection $I\times J\longrightarrow I$ est tjrs \deleted{\unsure{coh. fi.}} une $\mathbb D$-fibr. coh. \ldots\deleted{\unsure{fibr}} \deleted{\unsure{donc prenons comme}} Prenons p. ex. pour $Y'$
\diagram{\begin{tikzcd}[ampersand replacement=\&,row sep=small]
\& \unsure{a|o} \arrow[dl]\\
b \& \\
\& a' \arrow[ul]\\
\underbrace{\hphantom{bb}}_{\unsure{Y_1|Y'_1|X_1}} \& \underbrace{\hphantom{aa}}_{Y'_0}
\end{tikzcd}}
, somme amalgamée de deux copies de $Y=\Delta^1=$

alors $X'=X\times_Y Y'$ peut se voir comme somme amalgamée $X\amalg_{X_1}X$. Mais ça semble \uline{pas} être un contre-exemple. \deleted{\unsure{La question}} Il \unsure{faudra|faudrait} que $X_1\longrightarrow X'$ ne soit \uline{pas} une $\mathbb D$-équivalence
% archive sheet number: 71
\scanpage{73}
\mspage{79}

\uline{Proposition \unsure{2|3}} Soit $f:X\longrightarrow Y$ dans $\mathrm{Diag}$. Conditions équivalentes

a) Pour tout $Y''\xrightarrow{p}Y'\xrightarrow{\unsure{g|r}}Y$, d'où
\diagram{\begin{tikzcd}[ampersand replacement=\&]
X''\arrow[r,"q"]\arrow[d] \& X'\arrow[r]\arrow[d] \& X\arrow[d,"f"]\\
Y''\arrow[r,"p"] \arrow[ur,phantom,"\text{cart}" description] \& Y'\arrow[r]\arrow[ur,phantom,"\text{cart}" description] \& Y
\end{tikzcd}}
Si $p$ est une $\mathbb D$-équivalence, $q$ aussi.

\margin{$\longrightarrow$ il faut supposer $Y',Y''$ (Hot)-asphériques cf page \unsure{96|95}, et supposer vérifié l'axiome $\mathrm{Der}\,3\,\unsure{a'|a})$ de la page \unsure{97}.}

% slanted
\margin{\uuline{$b)\Rightarrow a)$ \unsure{canular}} cf annotation marginale cf contre-exemple page 89. Mais ça marche \unsure{p.} pour les dérivateurs $\mathbb D_{\unsure{i|\delta}}$ $(i\in\mathbb N)$ associés aux ``localisateurs fondamentaux'' $W_i$, i.e. la théorie des \unsure{$i$|$\delta$}-types d'homotopie}

% slanted
\margin{L'implication $b\Rightarrow a$ pas prouvée, seulement sous forme plus faible : a$'$) si $Y'',Y'$ \deleted{\unsure{sont}} $\mathbb D$-asphériques et 1-connexes, \deleted{\unsure{c'est alors que}} alors $X''\longrightarrow X'$ est une $\mathbb D$ équivalence. (Il suffit $=$ de \unsure{de le|le|la} supposer pour $Y'$, $Y''$ \deleted{\unsure{tel}} ayant des objets ``\unsure{finaux}'' (pour \unsure{avoir} qu'a$'$ $\Rightarrow$ b) ou ``ayant des objets initiaux'' pour impliquer \unsure{a$'$)$\Rightarrow$b$'$)}}

b) \deleted{\unsure{Donc}} $f$ est une $\mathbb D$-fibration cohomologique \unsure{univ\textsuperscript{lle}}, i.e. elle est une fibration \inserted{coh.}, et le \unsure{reste} par tout changement de base.

b$'$) Condition duale de b) : \deleted{\unsure{$\mathbb D$}} $f$ est une $\mathbb D$-fibration homologique universelle.

$a)\Longrightarrow b)$ Supposons a), Soit $Y'\to Y$ chgt de base, d'où
\diagram{\begin{tikzcd}[ampersand replacement=\&]
X'\arrow[r]\arrow[d,"f'"']\& X\arrow[d,"f"]\\
Y'\arrow[r]\&Y
\end{tikzcd}}
. Pour prouver que $f'$ est une $\mathbb D$-fib coh. on regarde au dessus de $Y'$ les foncteurs de la forme $\{y'\}\longrightarrow Y'_{/y'}$ $(y'\in\mathrm{Ob}\,Y')$ et $Y'_{/y'_1}\longrightarrow Y'_{/y'_0}$ (si $u':y'_1\unsure{-|\longrightarrow}y'_0$ flèche dans $Y'$), et on veut que les flèches \deleted{déduites} de $\mathrm{Cat}_{/Y'}$ déduites de celles-là par le chgt de base $X'\xrightarrow{f'}Y'$ soient des $\mathbb D$-équiv.
% archive sheet number: 72
\scanpage{74}
\mspage{80}
\diagram{\begin{tikzcd}[ampersand replacement=\&,column sep=small]
X^{\unsure{\prime|1}}\arrow[r,leftrightarrow,"q"]\arrow[d,"f'_1"'] \& X'_0\arrow[r,dash]\arrow[d,"f'_0"] \& X'\arrow[r,dash]\arrow[d,"f'"] \& X\arrow[d,dash,"f"]\\
Y'_1\arrow[r,leftrightarrow,"p"']\arrow[ur,phantom,"\text{cart}" description] \& Y'_0\arrow[r]\arrow[ur,phantom,"\text{cart}" description] \& Y'\arrow[r]\arrow[ur,phantom,"\text{cart}" description] \&Y
\end{tikzcd}}

Comme $p$ est une $\mathbb D$-équiv, cela résulte bien de a).

$b)\Longrightarrow a)$. Comme la condition b) est évidemment stable par chgt de base, on est \deleted{\unsure{réduit}} ramené à prouver que \deleted{$f$ lui-$=$ est une $\mathbb D$-fibration coh. Mais c'est évident} \deleted{sous la forme a) (prop. 1) de la condition en question. Comme a) est une condition \unsure{autoduale|auto-duale} (\unsure{ou} par raisonnement direct calqué sur le précédent)} si $f$ satisfait b), et si on a
\diagram{\begin{tikzcd}[ampersand replacement=\&]
\&X\arrow[d,"f"'] \& X'\arrow[l,"q"']\arrow[d,"f'"]\\
(*)\&Y \&Y'\arrow[l,"p"]\arrow[ul,phantom,"\text{cart}" description]
\end{tikzcd}}
si $p$ est une \inserted{$\mathbb D$-}équiv., alors $q$ aussi, i.e.
\[
\underset{\substack{\unsure{\shortparallel}\\H^\bullet_{\mathbb D}(Y,f_*(M_X))}}{H^{\unsure{\bullet|\cdot|\prime}}_{\mathbb D}(X,M)}
\mathrel{\mathop{\longrightarrow}\limits^\sim}
\underset{\substack{\unsure{\shortparallel}\\H^\bullet_{\unsure{\mathbb D|\cdot}}(Y',f'_*(M_{X'}))}}{H^\bullet_{\mathbb D}(X',M)}
\]
pour tt $M$ de $A=\mathbb D(e)$. Mais \deleted{\unsure{cette flèche}} $F=f_*(M_X)$ est loc. constant et se calcule ``fibre par fibre'', de $=$ pour $F'=f'_*(M_{X'})$.

% slanted
\margin{et on a donc $F'\simeq p^*F$. On est \deleted{\unsure{conduit}} ramené à admettre que \deleted{\unsure{$f$}} $p$ étant $\mathbb D$-équiv. \unsure{implique} l'iso. $H^\bullet_{\mathbb D}(Y,F)\longrightarrow H^\bullet_{\mathbb D}(Y,p^*F)$. \unsure{C'est sûrement|C'est évidemment} \deleted{\unsure{pour}} pour $F$ constant, mais est-ce vrai pour $F$ loc. constant ? Ce n'est pas évident.}

\deleted{\unsure{s'identifie}, grâce au fait que $f,f'$ sont des $\mathbb D$-fibr. coh, à \[
\deleted{\unsure{f}}\ H^\bullet_{\mathbb D}(Z,M)\longrightarrow H^\bullet_{\mathbb D}(Z',M)
\] \unsure{où} $Z$ est la fibre de $X$ en $y$, $Z'$ celle en $y'$, et la flèche est induite par $Z'\longrightarrow Z$. Comme le \unsure{carré} est cartésien, $Z'\longrightarrow Z$ est un iso, et on gagne.}
% archive sheet number: 73
\scanpage{75}
\mspage{81}

L'équiv. $a\Longleftrightarrow b')$ \deleted{\unsure{se}} se déduit de $a)\Longrightarrow b)$ par dualité (la condition a) étant \unsure{autoduale}), ou en reprenant le $\deleted{\unsure{\Leftrightarrow}}=$ raisonnement pour $a\Rightarrow b'$, et pour $b')\Rightarrow a)$ aussi (en \deleted{changeant} remplaçant la coh. par l'homologie).

\uline{Définition 2} \unsure{Si} \deleted{\unsure{de}} les conditions équivalentes de la prop. 2 sont satisfaites, on dit que $f$ est une \uline{$\mathbb D$-fibration}.

\margin{\uline{NB} C'est une notion à \uline{la fois} locale et ``colocale'' sur $Y$ (alors que fibr. g. est locale, fibr. dr. est colocale)}

\uline{Proposition 3} a) La notion de $\mathbb D$-fibration est stable par changt de base \uline{et par composition}\textsuperscript{*}

% slanted
\margin{?? Stabilité par composition \uline{n'est} pas claire}

b) Soient
\diagram{\begin{tikzcd}[ampersand replacement=\&]
X\arrow[rr,"f"]\arrow[dr,dash,"p"'] \& {} \& X'\arrow[dl,dash,"p'"]\\
\& Y\arrow[u,phantom,"\text{comm}" description] \&
\end{tikzcd}}
avec $p,p'$ des $\mathbb D$-fibrations \inserted{(il \unsure{suffira|suffirait} que $p$ et $p'$ soient des $\mathbb D$-fibr. coh., ou que ce soient des $\mathbb D$-fibr. homologiques, ou que $p,p'$ soient propres, ou (t.t.\unsure{lisses|lisse}))$;$} \deleted{\unsure{pour que}} $f$ induit des $\mathbb D$-équiv. sur les fibres \unsure{soit}; alors $f$ est une $\mathbb D$-équivalence.

% slanted
\margin{Peut se préciser considérablement dans le cas \deleted{\unsure{où}} où $p,p'$ sont des $\mathbb D$-fibrations. Cf prop. 6 page 84}

\uline{Dém} a) est immédiat sur la condition a) (prop 2 :) de la notion de $\mathbb D$-fibration. Pour b), supposons \inserted{que} $f$ induise des $\mathbb D$-équiv. sur les fibres. On veut prouver que
\[
H^{\unsure{\bullet|\cdot|\prime}}_{\mathbb D}(X,M)\unsure{-|\longrightarrow} H^\bullet_{\mathbb D}(X',M)
\]
est un iso, pour tt $M$ dans $A$, mais cette flèche s'identifie à
\[
H^\bullet_{\mathbb D}(Y,f_*(M_X))\longrightarrow H^\bullet_{\mathbb D}(Y,f'_*(M_{X'}))
\]
déduite de
\[
f_*(M_X)\unsure{-|\longrightarrow} f_*(M_{X'})\qquad(\mathrm{NB}\quad M_X\simeq f^*M_{X'}),
\]
% archive sheet number: 74
\scanpage{76}
\mspage{82}

et il suffit de prouver que cette flèche est un iso. Pour ceci, on regarde fibre par fibre, et on utilise le fait que ``\deleted{\unsure{les}} les images directes se calculent fibre par fibre''.

\uline{Question} Soit $f:X\longrightarrow Y$ une $\mathbb D$-fibration, et soit $F\in\mathrm{Ob}\,\mathbb D(Y)$. Alors $f_* f^*F$ se calcule-t-il \deleted{\unsure{fibre}} \unsure{p. r.|par} \deleted{\unsure{par|pro}} \inserted{aux} fibres ? Même question pour $f_! f^*F$. Pour la première propriété, concernant $f_*$, on montre que cela marche si $F$ est de la forme $p_*(\deleted{\unsure{N}}M_{Y'})$, où $p:Y'\longrightarrow Y$ est $\mathbb D$-\uline{propre}, et $M\in\mathrm{Ob}\,A$.

D'autre part, la sous-catégorie pleine \inserted{$C$} de $\mathbb D(Y)$ formée des $F$ pour lesquels ça marche est \inserted{(sauf erreur)} ``stable par $\mathbb D$-$\unsure{\mathrm{hot}|\mathrm{ho}}$-$\underset{\longleftarrow}{\mathrm{lim}}$'', en le sens suivant : si on a une catégorie $I$ avec objet final $e$ formant une sous-catégorie fermée (i.e. $e\longrightarrow x\ \Longrightarrow\ x=e$), et si on a $\mathscr F\in\deleted{\unsure{A}}\mathbb D(Y\times I)$, d'où
\[
\begin{aligned}
I^\circ&\longrightarrow\mathbb D(Y)\\
i&\longrightarrow\mathscr F_i
\end{aligned}
\]
(diagramme squelettique), si les $\mathscr F_i$ pour $i\ne e$ sont dans $C$, et si $\mathscr F$ est dans l'image essentielle de $j_*$, où $j:Y\times(I\setminus\{e\})\unsure{\longleftarrow|\longrightarrow}Y\times I$ est l'inclusion, alors $\mathscr F_e$ est aussi dans $C$.

On peut cependant considérer que si on veut que $f:X\longrightarrow Y$ satisfasse à toutes les propriétés coh. et hom. qui s'attachent à la notion d'une ``fibration'' (``loc\unsure{.} triviale''), on devrait mettre dans les conditions, que $f$ est \uline{lisse}, et \uline{propre}, donc que le chgt. de base par $f$ commute aux foncteurs $g_*,g_!$.
% archive sheet number: 75
\scanpage{77}
\mspage{83}

\deleted{Cela doit signifier aussi que la \unsure{convolution|construction|considération|conclusion} de} Signalons

\uline{Proposition} 4. Soit $f:X\longrightarrow Y$ \uline{propre}. Pour que $f$ soit une $\mathbb D$-fibration, il f. et s. que pour \deleted{tout} $y_1\longrightarrow y_0$ de $Y$, le foncteur $\deleted{\unsure{B}}\ X_{/y_1}\unsure{-|\longrightarrow}X_{/y_0}$ induit soit une $\mathbb D$-équiv. \inserted{(i.e que $f$ soit une $\mathbb D$-fibr. \uline{coh.})}. Dualement, si $f$ est $\mathbb D$-lisse, pour que ce soit une $\mathbb D$-fibration, il \inserted{f. et} suffit que $\forall\, y_1\longrightarrow y_0$ de $Y$, $ {}_{y_0\backslash}X\longrightarrow {}_{y_1\backslash}X$ soit une $\mathbb D$-équivalence.

\margin{\uline{Cor} \deleted{\unsure{$f$}} Supposons $f$ Cat-cofibrant (donc propre). Pour que $f$ soit une $\mathbb D$-fibration, il f. et s. que les foncteurs $X_{y_0}\longrightarrow X_{y_1}$ (pour $y_0\longrightarrow y_1$ de $Y$) soient des $\mathbb D$-équivalences. Dualement pour Cat-fibration (donc foncteurs lisses)}

On pourrait appeler $\mathbb D$-fibrations \uline{parfaites} les foncteurs à la fois propres et lisses qui sont des \inserted{$\mathbb D$-}fibr. coh. (ou ce qui revient au $=$, des $\mathbb D$-fibr. \unsure{droites|dro.}).

\deleted{Définition} \uline{Proposition} 5. Soit $f:X\longrightarrow Y$ \unsure{de|dans} $\mathrm{Diag}$. Conditions équivalentes

a) \deleted{\unsure{$\forall$}} $f$ est une $\mathbb D$ équiv., et le reste par tt changement de base

b) $f$ est une $\mathbb D$ fibr. à fibres \inserted{$\mathbb D$-}contractiles

\uline{Dém} $a\Longrightarrow b$ Que \deleted{\unsure{$f$}} a) implique que $f$ soit une $\mathbb D$-fibr. est immédiat sur la définition de celle-ci (condition a) de la prop. 2), et que les fibres sont contractiles est tautologique.

$b)\Longrightarrow a)$ Il suffit de montrer que b) implique que $f$ est une $\mathbb D$ équivalence, Mais c'est un cas particulier de Prop. 3 b) (en y prenant $X'=Y$).
% archive sheet number: 76
\scanpage{78}
\mspage{84}

Quand ces conditions sont satisfaites, on s'attend à ce que le foncteur
\[
f^*:\mathbb D(Y)\longrightarrow\mathbb D(X)
\]
soit pl. fid, i.e. que les morphismes d'adjonction
\[
F\unsure{-|\longrightarrow} \unsure{p_*|f_*|f} f^*F,\qquad f_! f^*\unsure{f|F}\unsure{-|\longrightarrow} F
\]
dans $\mathbb D(Y)$ soient des iso. (C'est vrai si $f$ est $\mathbb D$-propre ou $\mathbb D$-lisse, cf cor. 3 p.~71.) Cela revient encore à dire que les $p_*\unsure{f^*|p^*}F$ se calculent par passage aux fibres, \deleted{\unsure{et}} \inserted{ou encore} itou pour les $p_!p^*F$. Cela nous ramène à la question de la page \unsure{87|82|89},

\uline{Définition 3} On dit alors que $f$ est une $\mathbb D$ fibration \uline{triviale}.

\uline{Proposition} 6 (précise prop. 3 page 81) Considérons un triangle comm. dans $\mathrm{Diag}$
\diagram{\begin{tikzcd}[ampersand replacement=\&]
X\arrow[rr,"f"]\arrow[dr,"p"'] \& \& X'\arrow[dl,"p'"]\\
\& S\&
\end{tikzcd}}

a) \deleted{\unsure{b)}} Supposons que $p$ et $p'$ soient des $\mathbb D$-fibrations gauches, ou que $p$ et $p'$ soient des $\mathbb D$-fibrations droites. Soit \deleted{\unsure{$T$}} $S_0$ une partie de $S$ qui rencontre toute composante connexe de $S$. Pour que $f$ soit une $\mathbb D$-équivalence, il \uline{suffit} que $\forall s_0\in S_0$, $f$ induise une $\mathbb D$-équivalence $X_{s_0}\unsure{-|\longrightarrow} Y_{s_0}$.

% slanted
\margin{cf contre-exemple page 87}

b) Cette condition est aussi nécessaire \inserted{(?)} si $p,p'$ sont des $\mathbb D$-fibrations. Donc dans ce cas, $f$ est une $\mathbb D$-équivalence ssi $f$ induit des $\mathbb D$-équivalences fibre par fibre.

% slanted
\margin{? \uuline{Faux} Pas prouvé ! Même si $X'=S$ ce n'est pas clair, même $\mathbb D=\mathrm{HOT}_{\unsure{\mathrm{AB}|\mathrm{Ab}}}$, $f$ propre et lisse.}

% slanted
\margin{Mais OK. si $S$ $\unsure{\mathbb D\text{-contractile}|\text{contractile}}$\ldots}
% archive sheet number: 77
\scanpage{79}
\mspage{85}

Dém. a) Il suffit, en vertu de prop. 3, de montrer que l'hyp. implique que pour \uline{tout} $s\in S$, $f$ induit une $\mathbb D$-équivalence $X_s\xrightarrow{f_s}X'_s$. Cela résulte \deleted{\unsure{du le}} aussitôt du \deleted{lemme}

\uline{Lemme} \deleted{\unsure{sous les hyp}} Si $p$ et $p'$ sont tous deux \unsure{soient|soit} des $\mathbb D$-fibrations gauches, \unsure{soient|soit} des $\mathbb D$-fibr. droites, et si $s_1\longrightarrow s_0$ est une flèche dans $S$, alors $f_{s_1}$ est une $\mathbb D$-équiv. sss $f_{s_0}$ l'est.

(\deleted{Donc} \inserted{L'ens $S'\subset\mathrm{Ob}\,S$} des $s\in S$ tels que $f_s$ soit une $\mathbb D$-équiv., est une réunion de comp. connexes de $S$.)

Supposons p. ex. que $p,p'$ soient des $\mathbb D$-fibr. gauches. Considérons \inserted{pour $s\in S$} le diagramme
\diagram{\begin{tikzcd}[ampersand replacement=\&]
X_{s\deleted{\unsure{1}}}\arrow[r,"f_s"']\arrow[r,phantom,shift left=1.6ex,"{\deleted{\unsure{\longleftrightarrow|\longrightarrow}}}" description]\arrow[d] \& X'_{\unsure{s|s_0}}\arrow[d]\\
X_{/\unsure{s|s_0}}\arrow[r,"f_{/s}"] \& X'_{/\unsure{s|s_0}}
\end{tikzcd}}
Les flèches verticales sont des $\mathbb D$-équiv. par l'hyp. sur $p,p'$, donc $f_s$ est une \inserted{$\mathbb D$-}équiv. sss $f_{/s}$ l'est.

Considérons maintenant
\diagram{\begin{tikzcd}[ampersand replacement=\&]
X_{/s_1}\arrow[r,"f_{\unsure{s_1|/s_1}}"]\arrow[d] \& X'_{/s_1}\arrow[d]\\
X_{/s_0}\arrow[r,"f_{/s_0}"] \& X'_{/s_0}
\end{tikzcd}}
Ici encore, \unsure{les|l'} hyp sur $p,p'$ \unsure{impliquent|implique} que les flèches verticales sont des $\mathbb D$-équiv., donc $f_{s_1}$ est une $\mathbb D$-équiv. sss $f_{s_0}$ l'est.
% archive sheet number: 78
\scanpage{80}
\mspage{86}

\deleted{b) Il faut prouver que si $p,p'$ sont des $\mathbb D$-fibrations, et $f$ une $\mathbb D$-équiv.\deleted{\illegible{1}}, alors $f$ induit des $\mathbb D$-équiv. sur les fibres. Considérons, pour $M\in\mathrm{Ob}\,A$, l'hom induit par $f$}
\deleted{\[\underset{\substack{\shortparallel\\H'}}{\underline{p'_* (M_{X'})}}\longrightarrow\underset{\substack{\shortparallel\\H}}{\underline{p_* (M_X)}}\,.\]}

\deleted{\deleted{\unsure{l'hypothèse|L'hypothèse}} Comme $f$ est une $\mathbb D$-fibration, c'est un hom d'objets loc. constants \inserted{sur $S$}, dont les fibres sont les $H^{\unsure{\bullet|\cdot|\prime}}_{\mathbb D}(X',M)$, $H^\bullet_{\mathbb D}(X,M)$ respectivement. Il faut prouver que l'hom. est iso, et tout ce qu'on sait, c'est qu'il induit $H^\bullet_{\mathbb D}(S,H')\unsure{-|\to}H^\bullet_S(S,H)$ qui est un iso. Si $H,H'$ étaient des \deleted{\unsure{fai|fais}} objets \uline{constants} de $\mathbb D(S)$, soient $P_S,Q_S$, alors}

Si $Y'=S$, l'énoncé \inserted{non prouvé} dit que si $f\colon X\longrightarrow S$ est une $\mathbb D$-fibration et une $\mathbb D$-équivalence, alors ses fibres sont $\mathbb D$-contractiles. Ce n'est pas sûr même pour une projection $X=T\times S\longrightarrow S$. Si $\unsure{y|\eta|\gamma|\mathfrak y}=\mathrm{hot}_0(\unsure{S|T}),\ \unsure{\zeta|z}=\mathrm{hot}_0(T)$, la question se réduit à celle-ci : Si $\forall M$ dans $A$, \deleted{\unsure{on a}}, l'hom
\deleted{\[(\zeta\times y)\times_{\mathbb D} M\xrightarrow{\sim}y\times_{\mathbb D} M\quad\text{alors}\]}
\[\alpha_M\unsure{=|:}\deleted{\unsure{y\times M|\zeta\times M}}\ \zeta\times_{\mathbb D}M\longrightarrow M\]
satisfait
\[y\times_{\mathbb D}\alpha_M\colon\ y\times_{\mathbb D}(\zeta\times_{\mathbb D} M)\xrightarrow{\sim}y\times_{\mathbb D} M\quad\text{\uline{iso}},\]
alors $\alpha_M$ lui-même est iso.
% archive sheet number: 79
\scanpage{81}
\mspage{87}

Voici un \uuline{contre-exemple} que je donne \inserted{d'abord} dans le contexte des espaces topologiques, au lieu de $\mathrm{Cat}$. Je vais construire une fibration loc$^{\mathrm t}$ triviale et non triviale

% slanted
\margin{Voir un contre-exemple encore plus simple \unsure{pages|page} 93, 94. (Mais qui a l'inconvénient de ne pas être à fibres \unsure{connexes}.)}
\diagram{\begin{tikzcd}[ampersand replacement=\&,column sep=small]
X\arrow[d,"f"'] \&[-8pt] {\supset S^1}\\
S=S^1 \&
\end{tikzcd}}
de fibre $S^1$, telle que $f$ soit une $\mathbb D$-équivalence $\forall\,\mathbb D=\mathrm{Der}(k\text{-}\mathrm{Mod})$, $k$ étant un anneau où 2 est inversible.

% slanted
\margin{\uline{NB} Exemple beaucoup \deleted{\unsure{en|et}} plus simple encore : on prend $X=S^1$, et $X\to S$ le revêtement à 2 feuillets. Si on prend le rev. à \unsure{$n$|$-$} feuillets c'est \unsure{encore} OK si \unsure{$n$|$2$} est \unsure{inv.}}

\deleted{\unsure{Donc|Dan}} Cela signifie que
\[\deleted{\unsure{H|\mathbb Z}}H_{\unsure{\bullet|*}}(X,\deleted{\unsure{k}}\mathbb Z)\longrightarrow\deleted{\unsure{H|\mathbb Z}}H_{\bullet}(S,\deleted{\unsure{k}}\mathbb Z)\]
est un iso mod 2-torsion (noyau et conoyau sont annulés par une puissance de 2 -- en fait par 2 seul), ou encore $H_{\bullet}(X,k)\xrightarrow{\sim}H_{\bullet}(S,k)$ si 2 inv. dans $k$. Mais la fibre $S^1$ n'est pas $\mathbb D$-contractile (sauf si $k=0$ !)

\deleted{L'idée est de choisir la fibration telle} La fibration est donnée, à isom. \inserted{unique} près, en se fixant un point $s\in S$, et en se donnant \deleted{\unsure{alors|alor}} un autom. \inserted{$u$} de la fibre $X_s$ \deleted{\unsure{qui soit}} (``monodromie'') (la classe d'isom. du fibré ne dépend que de la classe d'isotopie \unsure{$u$|$\sim$}). Je prendrai pour $u$ la symétrie (renversant l'orientation) $z\longrightarrow z^{-1}$ du cercle \unsure{standart|standard} $\{z\in\mathbb C\mid |z|=1\}$.

Cette fibration ``se réalise'' \deleted{\unsure{par}} \deleted{\unsure{elle|de}} en un ordonné \inserted{$X\to S$}, $S^1$ étant représenté (\unsure{seule base|seule façon} \unsure{connue} sur les fibres) par
\diagram{\begin{tikzcd}[ampersand replacement=\&,column sep=small,row sep=small]
\& \& a \& \&\\
{S={}} \& \unsure{o|0|a}\arrow[ur]\arrow[dr] \& \& \unsure{o'|0'|a'}\arrow[ul]\arrow[dl] \& {=S'\amalg_R S''}\\
\& \& b \& \&
\end{tikzcd}}
où $R$ est la cat. discrète ayant deux éléments $\unsure{0|a},1$. On définit $X$ \deleted{\unsure{en}} comme $X'\amalg_Y X''$, où $X'=X\mid S'$, \deleted{\unsure{$X\mid S$}}, $X''=X\mid S''$, $Y=X\mid R$, \deleted{\unsure{avec $X=$|que $X=$}} $Y=X\mid R$.

% slanted
\margin{NB $S',S''$, \inserted{$R=S'\cap S''$} fermés dans $S$, $X',X''$, \inserted{$Y=X'\cap X''$} fermés dans $X$.}
% archive sheet number: 80
